\section{BUG-017: Git Clone Silence on Failure}
\label{sec:bug-017}

\subsection*{Metadata}
\begin{tabular}{ll}
\textbf{Issue ID} & BUG-017 \\
\textbf{Severity} & Medium \\
\textbf{Category} & Bug Fix \\
\textbf{Affected File} & \srcfile{src/loader.ts} \\
\textbf{Branch} & \branchref{fix/BUG-017-git-clone-silence} \\
\textbf{Changes} & \branchdiff{fix/BUG-017-git-clone-silence} \\
\end{tabular}

\subsection*{Executive Summary}
The \texttt{resolveInheritance} and \texttt{resolveDependencies} functions in \texttt{src/loader.ts} use \texttt{execSync} to run \texttt{git clone} commands with \texttt{2>/dev/null || true} appended. This shell construct redirects all stderr output from git to \texttt{/dev/null} and then forces the exit code to 0 (success) using \texttt{|| true}. The result is that \textbf{every git clone failure — regardless of cause — is completely invisible} to both the operator and the application logic.

When a parent agent repository is unreachable, the user agent directory has no \texttt{.gitconfig}, the source URL is malformed, the branch does not exist, or the disk is full, the error is swallowed before it reaches the try/catch block. The code then silently continues without the cloned dependency, and the agent loads with a degraded or misconfigured state.

The bug is at \texttt{src/loader.ts:180-183} (for inheritance) and at \texttt{src/loader.ts:226-229} (for dependencies), where \texttt{2>/dev/null || true} is appended to the \texttt{git clone} command inside the \texttt{execSync} call.

This is a diagnostic integrity issue: the errors exist but are discarded. When an agent operator sees unexpected behavior — missing parent rules, unresolved dependencies, or incomplete manifests — there is no error trail to investigate.

---

\subsection*{Root Cause Analysis}
\#\#\# The Code



\begin{lstlisting}[language=TypeScript]
// loader.ts:179-187 — resolveInheritance
try {
    execSync(`git clone --depth 1 "${manifest.extends}" "${parentDir}" 2>/dev/null || true`, {
        cwd: agentDir,
        stdio: "pipe",
    });
} catch {
    // Clone failed, continue without parent                            ← Dead code — catch is never reached
    return { manifest, parentRules: "" };
}

\end{lstlisting}





\begin{lstlisting}[language=TypeScript]
// loader.ts:223-233 — resolveDependencies
for (const dep of manifest.dependencies) {
    const depDir = join(depsDir, dep.name);
    try {
        execSync(
            `git clone --depth 1 --branch "${dep.version}" "${dep.source}" "${depDir}" 2>/dev/null || true`,
            { cwd: agentDir, stdio: "pipe" },
        );
    } catch {
        // Clone failed, skip this dependency                           ← Dead code — catch is never reached
    }
}

\end{lstlisting}



\#\#\# The Shell Redirection Anatomy

The command string \texttt{git clone --depth 1 "..." "..." 2>/dev/null || true} breaks down as follows:

| Component | Effect |
|---|---|
| \texttt{git clone --depth 1 "URL" "DIR"} | The actual clone operation |
| \texttt{2>/dev/null} | Redirects stderr (file descriptor 2) to \texttt{/dev/null} |
| \texttt{|| true} | If the previous command exits non-zero, run \texttt{true} (which exits 0) |

The \texttt{|| true} is the critical piece: it ensures that \texttt{execSync} itself never sees a non-zero exit code. Without \texttt{|| true}, \texttt{execSync} would throw an error when git clone fails. With \texttt{|| true}, \texttt{execSync} always returns successfully, and the try/catch block is \textbf{dead code} — it never executes.

\#\#\# The Two Failure Points

\textbf{Point 1: \texttt{resolveInheritance} (line 180)}

When \texttt{manifest.extends} points to a URL like \texttt{https://github.com/nonexistent/agent.git}:
1. \texttt{git clone} tries to connect — fails with \texttt{fatal: repository '...' not found}
2. Stderr with the error message goes to \texttt{/dev/null}
3. \texttt{|| true} converts the failure exit code to 0
4. \texttt{execSync} returns normally
5. The catch block is skipped
6. The code reads the cloned directory — which doesn't exist — and falls through to the inner catch (line 194)
7. \texttt{resolveInheritance} returns \texttt{\{ manifest, parentRules: "" \}} with no error

The operator sees nothing. The parent rules are silently omitted.

\textbf{Point 2: \texttt{resolveDependencies} (line 226)}

When a dependency source is unreachable:
1. \texttt{git clone} fails with a network error
2. Error is redirected to \texttt{/dev/null}, exit code neutralized
3. The catch block is dead code
4. The loop continues to the next dependency
5. No dependency files exist on disk
6. Any agent functionality that depends on the mounted dependency files silently breaks

\#\#\# Why This Is Medium Severity

1. \textbf{Complete diagnostic blindness}: Every git clone failure — network error, auth failure, missing repo, invalid URL — is invisible
2. \textbf{Dead code}: The try/catch blocks are misleading; the catch will never execute
3. \textbf{Silent degradation}: The agent loads with missing parent rules or missing dependencies but continues as if nothing is wrong
4. \textbf{Hard to debug}: Operators must inspect the actual \texttt{deps/} directory and compare with \texttt{agent.yaml} to detect missing clones
5. \textbf{Compounds other bugs}: BUG-015 (stale session state) and BUG-018 (duplicate dependency names) become harder to diagnose when the underlying git failure is hidden

---

\subsection*{Fix Implementation}
\#\#\# Option A: Remove \texttt{|| true} and Log Stderr (Recommended)

Remove the \texttt{|| true} construct to allow \texttt{execSync} to throw on failure, and log the actual git error:



\begin{lstlisting}[language=TypeScript]
// loader.ts — resolveInheritance
try {
    execSync(`git clone --depth 1 "${manifest.extends}" "${parentDir}" 2>&1`, {
        cwd: agentDir,
        stdio: "pipe",
        encoding: "utf-8",
    });
} catch (err: any) {
    const stderr = err.stderr || err.message || "unknown error";
    console.warn(`[loader] Failed to clone parent "${manifest.extends}": ${stderr.toString().trim()}`);
    return { manifest, parentRules: "" };
}

\end{lstlisting}


\begin{lstlisting}[language=TypeScript]
// loader.ts — resolveDependencies
for (const dep of manifest.dependencies) {
    const depDir = join(depsDir, dep.name);
    try {
        execSync(
            `git clone --depth 1 --branch "${dep.version}" "${dep.source}" "${depDir}" 2>&1`,
            { cwd: agentDir, stdio: "pipe", encoding: "utf-8" },
        );
    } catch (err: any) {
        const stderr = err.stderr || err.message || "unknown error";
        console.warn(`[loader] Failed to clone dependency "${dep.name}" from "${dep.source}": ${stderr.toString().trim()}`);
    }
}

\end{lstlisting}



\textbf{What this fixes:}
1. Git errors are captured and logged with \texttt{console.warn}
2. The try/catch blocks are no longer dead code
3. Operators can see exactly why a clone failed
4. The application continues gracefully (clone failures are still non-fatal)

\#\#\# Option B: Use \texttt{execFile} for Shell-Injection Safety

Replace \texttt{execSync} with \texttt{execFileSync} to avoid shell injection while preserving error output:



\begin{lstlisting}[language=TypeScript]
// loader.ts — resolveInheritance with execFileSync
import { execFileSync } from "child_process";

try {
    execFileSync("git", ["clone", "--depth", "1", manifest.extends, parentDir], {
        cwd: agentDir,
        stdio: "pipe",
        encoding: "utf-8",
    });
} catch (err: any) {
    const stderr = err.stderr || err.message || "unknown error";
    console.warn(`[loader] Failed to clone parent "${manifest.extends}": ${stderr.toString().trim()}`);
    return { manifest, parentRules: "" };
}

\end{lstlisting}



\textbf{Advantages:}
1. No shell injection risk (URL strings with special characters are safe)
2. Git errors are captured in stderr naturally
3. No need for \texttt{2>\&1} redirection

\textbf{Disadvantages:}
1. Slightly more verbose (array arguments)
2. \texttt{--branch} with empty version string requires conditionals

\#\#\# Fix Comparison

---

\subsection*{Affected Files}
\begin{itemize}
    \item \srcfile{src/loader.ts}
\end{itemize}

\subsection*{Verification}
The fix is verified through unit tests and integration tests that confirm the corrected behavior:

\begin{lstlisting}[language=TypeScript]
// verification/bug-017-verify.ts
import { describe, it, before, after } from "node:test";
import assert from "node:assert";
import { loadAgent } from "../src/loader";
import { rmSync, mkdirSync, writeFileSync, readFileSync } from "fs";
import { join } from "path";

const TEST_DIR = "/tmp/bug-017-verify";

function createAgent(yaml: string) {
    rmSync(TEST_DIR, { recursive: true, force: true });
    mkdirSync(TEST_DIR, { recursive: true });
    writeFileSync(join(TEST_DIR, "agent.yaml"), yaml.trim());
}

describe("BUG-017: Git Clone Silence on Failure", () => {
    after(() => {
        rmSync(TEST_DIR, { recursive: true, force: true });
    });

    it("should log a warning when inheritance clone fails", async () => {
        createAgent(`
name: "Test Agent"
version: "1.0.0"
description: "Parent clone test"
extends: "https://github.com/nonexistent-repo-12345/agent.git"
model:
  preferred: "anthropic:claude-sonnet-4-5-20250929"
  fallback: []
tools:
  - memory
        `);

        const warnings: string[] = [];
        const origWarn = console.warn;
        console.warn = (...args: any[]) => {
            warnings.push(args.join(" "));
            origWarn.apply(console, args);
        };

        await loadAgent(TEST_DIR);

        console.warn = origWarn;

        const hasCloneWarning = warnings.some(w => w.includes("clone") && w.includes("parent"));
        assert.ok(hasCloneWarning, "Should log a warning about the failed parent clone");
    });

    it("should log a warning when dependency clone fails", async () => {
        createAgent(`
name: "Test Agent"
version: "1.0.0"
description: "Dep clone test"
model:
  preferred: "anthropic:claude-sonnet-4-5-20250929"
  fallback: []
tools:
  - memory
dependencies:
  - name: "missing-dep"
    source: "https://github.com/nonexistent-repo-12345/missing.git"
    version: "main"
    mount: "/tmp"
        `);

        const warnings: string[] = [];
        const origWarn = console.warn;
        console.warn = (...args: any[]) => {
            warnings.push(args.join(" "));
            origWarn.apply(console, args);
        };

        await loadAgent(TEST_DIR);

        console.warn = origWarn;

        const hasDepWarning = warnings.some(w => w.includes("clone") && w.includes("dependency"));
        assert.ok(hasDepWarning, "Should log a warning about the failed dependency clone");
    });

    it("should still succeed when all clones work", async () => {
        // Create a valid agent.yaml without extends or deps
        createAgent(`
name: "Valid Agent"
version: "1.0.0"
description: "No deps needed"
model:
  preferred: "anthropic:claude-sonnet-4-5-20250929"
  fallback: []
tools:
  - memory
        `);

        const agent = await loadAgent(TEST_DIR);
        assert.strictEqual(agent.manifest.name, "Valid Agent", "Should load successfully");
    });

    it("should not suppress git error details in warning messages", async () => {
        createAgent(`
name: "Test Agent"
version: "1.0.0"
description: "Error detail test"
extends: "https://github.com/invalid-url-that-will-fail.git"
model:
  preferred: "anthropic:claude-sonnet-4-5-20250929"
  fallback: []
tools:
  - memory
        `);

        const warnings: string[] = [];
        const origWarn = console.warn;
        console.warn = (...args: any[]) => {
            warnings.push(args.join(" "));
            origWarn.apply(console, args);
        };

        await loadAgent(TEST_DIR);

        console.warn = origWarn;

        // The warning should contain git error details, not just a generic message
        const cloneWarning = warnings.find(w => w.includes("clone"));
        assert.ok(cloneWarning, "Should have a clone warning");
        // Should include either the URL or the error details
        assert.ok(
            cloneWarning!.includes("fatal:") ||
            cloneWarning!.includes("repository") ||
            cloneWarning!.includes("not found"),
            "Warning should include git error details"
        );
    });
});

\end{lstlisting}

